\begin{corollary}[F6: nonabelian pro-$p$ kernels over abelian tame quotients; Proved]
Under the equal-characteristic hypotheses of F5, every quotient $G$
of $\pi_1(U)$ which admits an exact sequence
\[
 1\longrightarrow P\longrightarrow G\longrightarrow\Gamma
 \longrightarrow1,
 \qquad P\text{ pro-}p,\quad
 \Gamma\text{ abelian pro-prime-to-}p,
\]
is topologically finitely generated. In fact one bound on the number
of generators works for all such quotients of the fixed $\pi_1(U)$.
The kernel $P$ is allowed to be nonabelian. The same finite-generation
conclusion holds when $\Gamma$ is a virtually abelian pro-prime-to-$p$
group, without asserting a uniform bound independent of its abelian
subgroup's index.
\end{corollary}
\begin{proof}
Let $d_0$ generate $\pi_1(U)^{\mathrm{ab},(p')}$, which is finitely
generated by F1. Then every $\Gamma$ in the first assertion has at
most $d_0$ generators. We bound $H^1(G,M)$ for every finite simple
$\mathbf F_\ell[G]$-module $M$.

For $\ell=p$, the normal pro-$p$ subgroup $P$ acts trivially on $M$.
Indeed its action factors through a finite $p$-group, whose every
nonzero characteristic-$p$ representation has a nonzero invariant
vector (induct on the group order using a central element of order
$p$ and the nonzero kernel of its nilpotent operator $z-1$).
The invariant subspace $M^P$ is $G$-stable, hence is all of the
simple module $M$. Thus $M$ has abelian prime-to-$p$ monodromy.
Inflation from $G$ to $\pi_1(U)$ is injective in degree one, and F5
therefore gives
\[
 \dim_{\mathbf F_p}H^1(G,M)\leq C_A\dim_{\mathbf F_p}M.
\]

For $\ell\ne p$, one has $H^1(P,M)=0$. A continuous cocycle factors
through a finite $p$-group together with its finite coefficient
action, and averaging by that group's invertible order makes it a
coboundary. Inflation--restriction now gives
\[
 H^1(G,M)\simeq H^1(\Gamma,M^P).
\]
Again $M^P$ is either $0$ or $M$ by simplicity and normality. In the
latter case evaluation on $d_0$ generators bounds the dimension of
this cohomology by $d_0\dim M$; in the former it is zero.

All finite quotients of $G$ are solvable, since they are extensions
of finite $p$-groups by finite abelian groups. Apply F4 with
$C=\max(C_A,d_0)$ to obtain at most $\lceil C\rceil+1$ generators,
uniformly in $G$.

If $\Gamma$ has an open abelian subgroup $\Gamma'$, its inverse
image $G'$ has finite index in $G$. The corresponding finite \'etale
cover $U'\to U$ has degree prime to $p$ and is the punctured spectrum
of a complete normal local surface $A'$ with the same algebraically
closed residue field. The first assertion applied to $A'$ shows
that $G'$ is finitely generated. Adding finitely many coset
representatives generates $G$.
\end{proof}

The class in F6 does not contain arbitrary solvable quotients:
alternating wild and tame normal factors can leave no normal
pro-$p$ subgroup with abelian tame quotient. Thus the full
prosolvable and full profinite problems remain separate gates.
